% 写作提示：按逻辑或时间顺序写过程，并保留关键异常的处理依据。
\section{分析、实验或实践过程}
\label{sec:process}

\subsection{实施条件}
示例只依赖已经记录的分项分数，不依赖特定软件、仪器或场地。需要重复计算时，可以把步骤写成算法或代码；课程不要求代码时，删去对应环境即可。

\subsection{核心过程}
\algoref{alg:score}是\eqref{eq:score}的逐步形式。算法标题放在上方，输入和输出单独列出。

\begin{algorithm}[htbp]
  \caption{由分项得分计算综合得分}
  \label{alg:score}
  \begin{algorithmic}[1]
    \Require 分项得分 $s_1,\ldots,s_n$ 与权重 $w_1,\ldots,w_n$
    \Ensure 综合得分 $S$
    \State $S \gets 0$
    \For{$i \gets 1,\ldots,n$}
      \State $S \gets S + w_i s_i$
    \EndFor
    \State \Return $S$
  \end{algorithmic}
\end{algorithm}

同一计算也可以写成代码，便于附上能够运行的实现。代码标题放在下方，见\listingref{lst:score}。

\begin{lstlisting}[language=Python,caption={按权重汇总分项得分},label={lst:score}]
def weighted_score(scores, weights):
    total = 0.0
    for score, weight in zip(scores, weights):
        total += score * weight
    return total
\end{lstlisting}

\subsection{关键问题及处理}
示例里没有缺失值。若某一项没有记录，应先决定是补计、剔除，还是单独报告，并在\tabref{tab:scores}的表注里写明。
